\documentclass[a4paper,12pt]{article}
%%%%%%%%%%%%%%%%%%%%%%%%%%%%%%%%%%%%%%%%%%%%%%%%%%%%%%%%%%%%%%%%%%%%%%%%%%%%%%%%%%%%%%%%%%%%%%%%%%%%%%%%%%%%%%%%%%%%%%%%%%%%%%%%%%%%%%%%%%%%%%%%%%%%%%%%%%%%%%%%%%%%%%%%%%%%%%%%%%%%%%%%%%%%%%%%%%%%%%%%%%%%%%%%%%%%%%%%%%%%%%%%%%%%%%%%%%%%%%%%%%%%%%%%%%%%
\usepackage{eurosym}
\usepackage{vmargin}
\usepackage{amsmath}
\usepackage{graphics}
\usepackage{framed}
\usepackage{epsfig}
\usepackage{subfigure}
\usepackage{fancyhdr}

\setcounter{MaxMatrixCols}{10}
%TCIDATA{OutputFilter=LATEX.DLL}
%TCIDATA{Version=5.00.0.2570}
%TCIDATA{<META NAME="SaveForMode"CONTENT="1">}
%TCIDATA{LastRevised=Wednesday, February 23, 201113:24:34}
%TCIDATA{<META NAME="GraphicsSave" CONTENT="32">}
%TCIDATA{Language=American English}

\pagestyle{fancy}
\setmarginsrb{20mm}{0mm}{20mm}{25mm}{12mm}{11mm}{0mm}{11mm}
\lhead{MA4603} \rhead{Kevin O'Brien} \chead{Midterm
Assessment Paper - Version A } %\input{tcilatex}

\begin{document}
\begin{center}
	\includegraphics[scale=0.60]{images/shieldtransparent2}
\end{center}

\begin{center}
	\vspace{1cm}
	\large \bf {FACULTY OF SCIENCE AND ENGINEERING} \\[0.5cm]
	\normalsize DEPARTMENT OF MATHEMATICS AND STATISTICS \\[1.25cm]
	\large \bf {MID-TERM ASSESSMENT EXAMINATION 1} \\[1.5cm]
\end{center}

\begin{tabular}{ll}
	MODULE CODE: MA4603 & SEMESTER: Autumn 2016\\[1cm]
	MODULE TITLE: Science Mathematics 3  & DURATION OF EXAM: 45 minutes \\[1cm]
	LECTURER: Mr. Kevin O'Brien & GRADING SCHEME: 20 marks \\
	& \phantom{GRadiC} \footnotesize {20\% of total module marks} \\[0.2cm]
	\\[1cm]
\end{tabular}
\begin{center}
	{\bf INSTRUCTIONS TO CANDIDATES}
\end{center}
\begin{itemize} 
	\item This exam will start at 12:05, and will last 45 minutes.
	
	\item Each question will be worth either 1 or 2 Marks. There are 15 Marks worth of questions.
	\item All questions must be attempted (LENS students please see below)
	
	\item Write all of your answers in the exam script. Write the script number on any other documents you submit.
	
	\item It is your responsibility to return the script to collection box. An audit of scripts will take place immediately after the exam. If your script is account for in that audit,  you are deemed to be absent, and will receive no marks.
	
	\item \textbf{IMPORTANT for LENS Student:}
	Specifically approved LENS students have to answer any selection of questions that have an aggregate mark of 12 Marks.  
	\begin{itemize}
		\item They may skip any three of the 1-Mark Questions
		\item OR - They may skip a 1-Mark Question and a 2-Mark Question
		\item The mark will be rescaled by 125 \%.
		\item They are advised to skip questions that are indicated by an asterisk symbol (``$\ast$"), but it is not compulsory that they do so.
	\end{itemize}
	
	
\end{itemize}
\newpage
\section*{Attempt ALL questions}

\subsection*{Q1. Dixon Q Test For Outliers (4 Marks)}

The typing speeds for one group of 12 Engineering students were recorded both at the beginning of year 1 of their studies. The results (in words per minute) are given below:

\begin{center}
	\begin{tabular}{|c|c|c|c|c|c|}
		\hline
		% Subject& A& B& C& D& E &F &G &H \\ \hline
		118 & 146 & 149 & 142 & 170& 153\\ \hline
		137 & 161 & 156& 165&  178& 159
		\\ \hline
	\end{tabular}
\end{center}
Use the Dixon Q-test to determine if the lowest value (118) is an outlier. You may assume a significance level of 5\%.
\begin{itemize}
	\item[i.](1 Mark)	State the Null and Alternative Hypothesis for this test.
	\item[ii.](1 Marks) Compute the test statistic
	\item[iii.](1 Mark) State the appropriate critical value.
	\item[iv.](1 Mark) What is your conclusion to this procedure
\end{itemize}



\subsubsection*{Part B} %NORMAL %8 MARKS
Assume that the length of injected moulded plastic components are normally distributed with a mean of 12.5mm and a standard deviation of 2.5mm.  Calculate the corresponding probability for the following measurements occurring on an individual component.

\begin{itemize}
	\item [i.](2 Marks)	Between 12.5 and 15mms,
	\item [ii.](2 Marks) Less than 10 mms,
	\item [iii.](2 Marks) Between 12 and 15 mms,
	\item [iv.](2 Marks) Less than 10.3 mms.
\end{itemize}
\noindent Use the normal tables to determine the probabilities for the above exercises. You are required to show all of your workings.
\newpage
(Write Your Answers Here)
\newpage
\vspace{0.25cm}
\subsubsection*{Part C} %10 MARKS
Data on the durations (measured in months) were collected for a random sample of product development projects.
The durations for these development projects were collected and tabulated as follows:

\begin{table}[ht]
	\begin{center}
		\begin{tabular}{|rrrrrrrr|}
			
			\hline
			12 & 11 & 20 & 19 & 18 & 9 & 16 & 15 \\
			\hline
		\end{tabular}
	\end{center}
\end{table}

\begin{itemize}
	\item[i.](1 Mark) Calculate the mean of the project durations.
	\item[ii.](2 Marks) Calculate the variance for this sample.
	\item[iii.](1 Mark) Calculate the standard deviation for this sample.
\end{itemize}

\end{document}